\documentclass[10.5pt]{article}
\usepackage[T1]{fontenc}
\usepackage[utf8]{inputenc}


\usepackage[margin=0.7in]{geometry}
\usepackage{enumitem}
\usepackage{titlesec}
\usepackage[hidelinks,colorlinks=true,urlcolor={darkred},linkcolor={darkred}]{hyperref}
\usepackage{xcolor}
\definecolor{darkred}{HTML}{8B0000}
\usepackage{parskip}
\setlength{\parskip}{4pt}

\newcommand{\headingfont}{\sffamily}
\titleformat{\section}{\headingfont\large\bfseries}{}{0em}{}[\vspace{-6pt}\rule{\textwidth}{0.6pt}]
\titlespacing*{\section}{0pt}{8pt}{4pt}
\pagestyle{empty}

\usepackage{fontawesome5}
\newcommand{\nameblock}{%
	{\headingfont\LARGE\bfseries Amin Shirazian}\\[2pt]
	\faMapMarker*\ Toronto, ON \;\textbar\; \faPhone\ +1 (647) 686-9324 \;\textbar\; \href{mailto:amin.shirazian.work@gmail.com}{\faEnvelope\ amin.shirazian.work@gmail.com} \;\textbar\; \href{https://aminshz.github.io/my-website/}{\faGlobe\ Website} \;\textbar\; \href{https://linkedin.com/in/amin-shirazian}{\faLinkedin\ LinkedIn}\\[-4pt]
	\rule{\textwidth}{1pt}}

	\newcommand{\entry}[2]{%
  \textbf{#1}\hfill{#2}\par\vspace{2pt}%
}
\begin{document}
	
	%======================= PAGE 1: TEACHING RESUME =======================
	\nameblock
	
	\vspace{2pt}
	{\headingfont\bfseries Economics Lecturer} --- six university course sections taught as instructor of record (2024--2026), including adult and continuing-education learners through TMU's Chang School; 13 graduate and undergraduate TA appointments; PhD in Economics expected Summer 2026.
	
	\section{Education}

	\entry{Ph.D., Economics --- Toronto Metropolitan University --- GPA 4.17/4.33}{August 2026}
	
	\vspace{4pt}
	\entry{M.A., Economics --- Tehran Institute for Advanced Studies --- GPA 4/4.33}{2020}
	\entry{B.A., Economics --- University of Tehran}{2018}
	
	% ===================== TEACHING =====================

	\section{Teaching Experience -- Toronto Metropolitan University}
	\textit{Complete list of courses, terms, and section sizes:} \href{https://aminshz.github.io/my-website/teaching.html}{aminshz.github.io/my-website/teaching.html}

	\subsection*{Contract Lecturer (2024--2026)}
	
	\begin{itemize}[leftmargin=*, itemsep=1pt, topsep=2pt]
		\item CECN603 Economic Issues in Globalization --- Spring 2026 (53 students); Summer 2025 (44)
		\item CECN301 Intermediate Macroeconomics I --- Fall 2025 (38); Winter 2025 (61)
		\item ECN301 Intermediate Macroeconomics I ---  Winter 2026 (87); 
		\item CECN204 Introductory Macroeconomics --- Summer 2024 (59)

	\end{itemize}

	Duties: full responsibility for syllabus delivery, lectures, assessment design, grading, and student support. 

	\subsection*{Graduate Teaching Assistant (2021--2026)}
	\begin{itemize}[leftmargin=*, itemsep=1pt, topsep=2pt]
		\item PhD level: Advanced Macroeconomics I (3 terms); Advanced Microeconomics II (2 terms); Mathematical Economics
		\item MA level: Macroeconomics (2 terms)
		\item Undergraduate: Intermediate \& Introductory Macro, Game Theory, Money and Banking, Economics of Information, Applied Economic Analysis, Math for Economics I \& II
	\end{itemize}
	
	Duties: tutorials, office hours, grading, and feedback.

		
	% ===================== RESEARCH =====================
	\section{Research \& Currency}

	\begin{itemize}[leftmargin=*, itemsep=1pt, topsep=2pt]
	\item \href{https://aminshz.github.io/my-website/paper1.pdf}{\textbf{Marginal Utility Shocks and the Precautionary Saving Puzzle}} \textit{(Job Market Paper)} ---
		Develops and calibrates a structural model linking health, longevity, and saving decisions to the HRS;
		informs retirement security and late-life wealth distribution.

	\item \href{https://aminshz.github.io/my-website/paper2.pdf}{\textbf{Means-Tested Insurance and Pareto-Tailed Wealth Distribution}} \textit{(working paper)} ---
		Imputes total medical spending in the HRS using MEPS; provides formal results on estimator consistency
		and documents a sharp dip in spending among middle-wealth households.

	\end{itemize}


	\section{Awards}
	Graduate Assistant Recognition Award, Toronto Metropolitan University --- 2022 and 2023 (awarded for excellence in teaching support and student service)

	
	\section{Technical \& Online-Delivery Skills}
	Zoom, D2L Brightspace, online assessment design \;\textbar\; Python, R, Stata, MATLAB, C++, \LaTeX \;\textbar\; data visualization and live coding demonstrations in class
	
	\newpage
	%======================= PAGE 2: COVER LETTER =======================
	\nameblock
	\begin{center}{\headingfont\Large\bfseries Cover Letter}\end{center}	
	
	\vspace{6pt}
	\today
	
	% \textcolor{red}{[Hiring Committee]}\\
	% \textcolor{red}{[Department / School]}\\
	% \textcolor{red}{[Institution]}
	
	Dear Members of the Hiring Committee,
	
	
	CECN 230 Math Econ II M Classroom Winter 2027 - Unit 2 (Job ID = 395940)

	% Course Code, Course Name, Number of Sections, Term(s) and Year
	% OR
	% Package X, Course Codes, Course Names, Number of Sections of Each Course, Term(s) and Year

	I am writing to apply for the position of {CECN 230 Math Econ II M Classroom Winter 2027 - Unit 2 (395940)}. I am completing my PhD in Economics at Toronto Metropolitan University (expected August 2026), and over the past two years I have taught six university course sections as instructor of record while serving as a teaching assistant in thirteen graduate and undergraduate courses. I believe my record demonstrates exactly the combination your posting asks for: current disciplinary knowledge, proven effectiveness with adult and diverse learners, and strong online teaching capability.
	
	My knowledge is current because I am an active researcher. My dissertation builds a quantitative model of household saving under health risk, estimated on recent U.S. survey data, and my lectures draw directly on this research practice: when I teach intermediate macroeconomics. I also regularly update my knowledge through checking new mainstream courses at different schools. 
	
	My lecturing experience is largely with adult learners. My CECN courses run through TMU's Chang School of Continuing Education, whose students are predominantly working adults returning to study part-time. Teaching this population assures that I have experience not only in theory, but in practice.
	
	These same Chang School courses are delivered online, so my online teaching is tested at scale: live Zoom lectures, course management on D2L Brightspace, online assessment design, and active discussion boards. Across my TA and lecturer roles I have handled student inquiries by email, and my Graduate Assistant Recognition Awards (2022, 2023) reflect that responsiveness.   
	
	I would welcome the chance to discuss how I can contribute to your students' success. My teaching statement, diversity statement, and full materials are enclosed; my website (\href{https://aminshz.github.io/my-website/}{aminshz.github.io/my-website}) hosts course details and references.\\
	
	Sincerely,\\[4pt]
	Amin Shirazian
	
	\newpage
	%======================= PAGE 3: TEACHING STATEMENT =======================
	\nameblock
	\begin{center}{\headingfont\Large\bfseries Statement of Teaching Philosophy}\end{center}
	
	Economics is a toolkit before it is a body of facts, and my central goal as a teacher is that students leave my course able to use the tools on problems I never showed them. My approach is essentialist in foundation and constructivist in delivery: I insist on rigorous command of the core analytical methods --- supply and demand, equilibrium, optimization, the aggregate models of modern macroeconomics --- but I teach them through cases, current events, and structured discussion rather than as formalism for its own sake.
	
	\textbf{Keeping content current.} Each term I teach Economic Issues in Globalization, I read a large number of current trade articles that students propose for their term paper, and I review their reports analyzing those articles. This keeps my own knowledge of trade issues up to date. For the paper, students have to find a recent article that makes a clear argument, take a position on it, and defend that position using data and the models from the course. In teaching Intermediate Macroeconomics, I rely on platforms that connect the models to current news and are updated on a regular basis. Staying current with these updates pushes me to continually refresh my own grasp of recent economic events and data.

	\textbf{Equity, Diversity, and Inclusion.}  I try to make my classroom open: I tell students from the first lecture that they are welcome to interrupt me, ask questions, or correct me if they think I am wrong, and I follow up over email and office hours. Because my students come from many different backgrounds, I pay particular attention to language barriers. I slow my pace, write down what I am saying as I go, and provide transcripts and captions so that students who are still building their English can follow the material fully.For students who face other challenges, I point them to Academic Accommodation Support early, invite them to raise any barriers at the start of the course, and offer flexible deadlines and academic consideration when circumstances require it.
	
	\textbf{Fairness.} I care deeply about fairness, and I believe an instructor's personal feelings about a student should never influence how that student is graded. Grades must reflect the quality of the work and nothing else. I hold myself to this standard closely, and I think it is one of the clearest ways an instructor can show respect for every student.

	\textbf{Accommodation} My philosophy is that life is bigger than any single course. I try to be very accommodating toward my students' needs and circumstances. I stay in regular contact with them by email, listen to their concerns, and do what I can to support them through difficulties so that a hard situation outside class does not derail their progress in it.

	\textbf{Engagement} I work to keep students connected to the course on several fronts. I stay in close contact with them I reply to individual concerns by email in detail, meet one-on-one over Zoom during office hours, and post regular announcements and reminders on D2L so no one falls behind. In class, I invite students to interrupt me and join the discussion rather than sit through a lecture, because I find they learn more when they are part of the conversation.\\
	
	Sincerely,\\[4pt]
	Amin Shirazian

	   
\end{document}